\chapter{The Bochner Technique}
\label{chap:pet-the-bochner-technique}

The Bochner identities relate connection Laplacians to curvature terms.  This
chapter records the resulting vanishing, dimension, and rigidity statements
for forms and other natural tensor bundles.

\section{Hodge Theory}

\begin{definition}[de Rham complex and cohomology]
  \label{def:pet-ch9-de-rham-cohomology}
  \lean{PetersenLib.deRhamCohomology}
  \source{petersen:page-0348}
  For a manifold $M$, let $\Omega^k(M)$ denote its smooth $k$-forms and let
  $d^k:\Omega^k(M)\to\Omega^{k+1}(M)$ be exterior differentiation.  The
  resulting \emph{de Rham complex} is
  \[
    0\to\Omega^0(M)\xrightarrow{d^0}\Omega^1(M)\xrightarrow{d^1}\Omega^2(M)\to\cdots
    \xrightarrow{d^{n-1}}\Omega^n(M)\to0,
  \]
  Its cohomology groups are
  $H^k(M)=\ker d^k/\operatorname{im}d^{k-1}$ and agree with the real
  cohomology of $M$.  In particular, $H^0(M)\cong\mathbb R$ for connected
  $M$, while $H^n(M)\cong\mathbb R$ when $M$ is compact and orientable.
\end{definition}

\begin{remark}[Poincaré duality pairing]
  \label{rem:pet-ch9-poincare-duality}
  \lean{PetersenLib.poincareDualityPairing}
  \source{petersen:page-0348,petersen:page-0349}
  \uses{def:pet-ch9-de-rham-cohomology}
  On a compact oriented $n$-manifold, the map
  \[
    (\omega_1,\omega_2)\longmapsto\int_M\omega_1\wedge\omega_2
  \]
  descends from $\Omega^k(M)\times\Omega^{n-k}(M)$ to a nondegenerate pairing
  $H^k(M)\times H^{n-k}(M)\to\mathbb R$.  Consequently these two cohomology
  groups are dual and have equal dimension.
\end{remark}

\begin{definition}[codifferential and normalized tensor inner product]
  \label{def:pet-ch9-codifferential}
  \lean{PetersenLib.codifferential,PetersenLib.normalizedTensorInnerProduct}
  \source{petersen:page-0349}
  \uses{def:pet-ch9-de-rham-cohomology}
  On a Riemannian manifold, the \emph{codifferential}
  $\delta^k:\Omega^{k+1}(M)\to\Omega^k(M)$ is the adjoint $\nabla^*$ of $d^k$;
  equivalently,
  \[
    \int_Mg(\delta^k\omega_1,\omega_2)\,\mathrm{vol}=\int_Mg(\omega_1,d^k\omega_2)\,\mathrm{vol}.
  \]
  For fixed-type tensors, use the normalized pairing
  \[
    (T_1,T_2)=\frac1{\operatorname{vol}M}
      \int_M g(T_1,T_2)\,\mathrm{vol}.
  \]
  The operators $d$ and $\delta$ remain adjoint for this normalization.
\end{definition}

\begin{definition}[Hodge Laplacian]
  \label{def:pet-ch9-hodge-laplacian}
  \lean{PetersenLib.hodgeLaplacian}
  \source{petersen:page-0349}
  \uses{def:pet-ch9-codifferential}
  Define the \emph{Hodge Laplacian} on $k$-forms by
  \[
    \square=d\delta+\delta d:\Omega^k(M)\longrightarrow\Omega^k(M).
  \]
  With the present signs, its restriction to functions equals $-\Delta$.
\end{definition}

\begin{lemma}[harmonic forms are closed and coclosed]
  \label{lem:pet-ch9-harmonic-iff-closed-coclosed}
  \lean{PetersenLib.hodgeLaplacian_eq_zero_iff}
  \source{petersen:page-0349}
  \uses{def:pet-ch9-hodge-laplacian}
  \dcref{ch9:9.1.1}
  A form $\omega$ satisfies $\square\omega=0$ precisely when
  $d\omega=\delta\omega=0$.  An exact form with this property is necessarily
  zero.
\end{lemma}
\begin{proof}
  \uses{lem:pet-ch9-harmonic-iff-closed-coclosed}
  \sketch
  Adjointness gives
  \[
    (\square\omega,\omega)=\|\delta\omega\|_2^2+\|d\omega\|_2^2.
  \]
  Thus a harmonic form is closed and coclosed, and the converse follows from
  the definition of \(\square\).  If also \(\omega=d\theta\), then
  \(\|\omega\|_2^2=(\theta,\delta\omega)=0\).
\end{proof}

\begin{definition}[Hodge cohomology]
  \label{def:pet-ch9-hodge-cohomology}
  \lean{PetersenLib.hodgeCohomology}
  \source{petersen:page-0350}
  \uses{def:pet-ch9-hodge-laplacian}
  The space of harmonic $k$-forms, called the \emph{Hodge cohomology}, is
  \[
    \mathcal H^k(M)=\ker(\square|_{\Omega^k(M)})
      =\{\omega\in\Omega^k(M):\square\omega=0\}.
  \]
\end{definition}

\begin{theorem}[harmonic representatives for de Rham cohomology]
  \label{thm:pet-ch9-hodge-theorem}
  \lean{PetersenLib.hodgeCohomology_equiv_deRhamCohomology}
  \source{petersen:page-0350}
  \uses{def:pet-ch9-hodge-cohomology, lem:pet-ch9-harmonic-iff-closed-coclosed, def:pet-ch9-de-rham-cohomology}
  Every de Rham class on a closed oriented Riemannian manifold has exactly one
  harmonic representative; equivalently,
  $\mathcal H^k(M)\xrightarrow{\sim}H^k(M)$.
\end{theorem}
\begin{proof}
  \uses{thm:pet-ch9-hodge-theorem,rem:pet-ch9-exercise-9-6-4,rem:pet-ch9-exercise-9-6-5}
  \sketch
  The map is well defined and injective by
  \cref{lem:pet-ch9-harmonic-iff-closed-coclosed}.  Elliptic spectral theory,
  developed in \cref{rem:pet-ch9-exercise-9-6-4} and
  \cref{rem:pet-ch9-exercise-9-6-5}, yields
  \(\Omega^k=\operatorname{im}\square\oplus\ker\square\).  Write a closed form as
  \(\omega=d\delta\theta+\delta d\theta+\widetilde\omega\), with
  \(\square\widetilde\omega=0\).  Applying \(d\) shows \(d\theta\) is harmonic;
  \cref{lem:pet-ch9-harmonic-iff-closed-coclosed} then eliminates
  \(\delta d\theta\).  Hence \(\omega\) and \(\widetilde\omega\) are cohomologous.
\end{proof}

\section{1-Forms}

Hodge theory reduces the computation of the first Betti number
$b_1(M)=\dim\mathcal{H}^1(M)$ to the study of harmonic $1$-forms; the same
program is generalized to other forms and tensors in later sections.

\subsection{The Bochner Formula}

For a harmonic $1$-form $\theta$ on $(M,g)$ let $X$ be the dual vector field,
$\theta(v)=g(X,v)$, and $f=\tfrac12|\theta|^2=\tfrac12|X|^2=\tfrac12\theta(X)$.

\begin{proposition}[duality between closed forms and symmetric derivatives]
  \label{prop:pet-ch9-vector-field-1form-duality}
  \lean{PetersenLib.dualVectorField_symmetric_iff_closed}
  \source{petersen:page-0350,petersen:page-0351}
  \uses{def:pet-ch9-de-rham-cohomology, def:pet-ch9-codifferential}
  \dcref{ch9:9.2.1}
  Let $\theta=X^\flat$, so $\theta(v)=g(v,X)$.  Then
  \begin{enumerate}
    \item the endomorphism $v\mapsto\nabla_vX$ is self-adjoint exactly when
      $\theta$ is closed;
    \item $\delta\theta=-\operatorname{div}X$.
  \end{enumerate}
\end{proposition}
\begin{proof}
  \uses{prop:pet-ch9-vector-field-1form-duality}
  \sketch
  The identity
  \(d\theta(V,W)+(\mathcal L_Xg)(V,W)=2g(\nabla_VX,W)\) separates the
  alternating and symmetric parts of \(\nabla X\), proving the first claim.
  Tracing the covariant derivative of the metric-dual relation gives
  \(\delta\theta=-\operatorname{div}X\).
\end{proof}

Consequently, when $\theta$ is harmonic then $\operatorname{div}X=0$ and $\nabla X$ is a
symmetric $(1,1)$-tensor; the Bochner formula for closed $1$-forms is stated
below through the dual vector field.

\begin{proposition}[Bochner formula for closed $1$-forms]
  \label{prop:pet-ch9-bochner-formula-1forms}
  \lean{PetersenLib.bochnerFormula_oneForm}
  \source{petersen:page-0351}
  \uses{prop:pet-ch9-vector-field-1form-duality}
  \dcref{ch9:9.2.2}
  Suppose $\nabla X$ is symmetric, or equivalently $X^\flat$ is closed.  Put
  $f=\tfrac12|X|^2$ and choose a local potential $X=\nabla u$ near $p$.  At
  $p$ one has
  \begin{enumerate}
    \item $\nabla f=\nabla_XX$.
    \item $\operatorname{Hess}f(V,V)=\operatorname{Hess}^2u(V,V)+(\nabla_X\operatorname{Hess}u)(V,V)+R(V,X,X,V)
      =|\nabla_VX|^2+g(\nabla^2_XX,V)+R(V,X,X,V)$.
    \item $\Delta f=|\operatorname{Hess}u|^2+D_X\Delta u+\operatorname{Ric}(X,X)
      =|\nabla X|^2+D_X\operatorname{div}X+\operatorname{Ric}(X,X)$.
  \end{enumerate}
\end{proposition}
\begin{proof}
  \uses{prop:pet-ch9-bochner-formula-1forms}
  \sketch
  Symmetry of \(\nabla X\) gives
  \(g(\nabla f,V)=g(\nabla_VX,X)=g(\nabla_XX,V)\).  Differentiate this identity
  once more, commute the two covariant derivatives, and use
  \(X=\nabla u\) to obtain the Hessian formula.  Tracing in a frame parallel
  along \(X\) at the point of evaluation gives
  \(\Delta f=|\nabla X|^2+D_X\operatorname{div}X+\operatorname{Ric}(X,X)\).
\end{proof}

\subsection{The Vanishing Theorem}

\begin{theorem}[Bochner vanishing for harmonic $1$-forms]
  \label{thm:pet-ch9-bochner-vanishing-1forms}
  \lean{PetersenLib.bochner_harmonicOneForm_parallel}
  \source{petersen:page-0352}
  \uses{prop:pet-ch9-bochner-formula-1forms}
  \dcref{ch9:9.2.3}
  Every harmonic one-form on a compact Riemannian manifold with
  $\operatorname{Ric}\ge0$ is parallel.
\end{theorem}
\begin{proof}
  \uses{thm:pet-ch9-bochner-vanishing-1forms}
  \sketch
  For the dual field \(X\), harmonicity and
  \cref{prop:pet-ch9-bochner-formula-1forms} give
  \[
    \Delta\tfrac12|X|^2=|\nabla X|^2+\operatorname{Ric}(X,X)\ge0.
  \]
  Compactness and the maximum principle make \(|X|\) constant; both
  nonnegative terms therefore vanish, so \(X\) and its dual form are parallel.
\end{proof}

\begin{corollary}[vanishing under pointwise positive Ricci curvature]
  \label{cor:pet-ch9-harmonic-1forms-vanish-positive-ricci}
  \lean{PetersenLib.harmonicOneForms_vanish_of_positiveRicci}
  \source{petersen:page-0352}
  \uses{thm:pet-ch9-bochner-vanishing-1forms}
  \dcref{ch9:9.2.4}
  Under the assumptions of \cref{thm:pet-ch9-bochner-vanishing-1forms}, if
  $\operatorname{Ric}$ is positive definite at some point, then
  $\mathcal H^1(M)=0$.
\end{corollary}
\begin{proof}
  \uses{cor:pet-ch9-harmonic-1forms-vanish-positive-ricci}
  \sketch
  By \cref{thm:pet-ch9-bochner-vanishing-1forms}, the dual field \(X\) is
  parallel and \(\operatorname{Ric}(X,X)=0\).  Positive definiteness at one
  point forces \(X\) to vanish there, hence everywhere.
\end{proof}

\begin{corollary}[first Betti bound and flat-torus rigidity]
  \label{cor:pet-ch9-betti1-bound-flat-torus}
  \lean{PetersenLib.firstBettiNumber_le_dim_iff_flatTorus}
  \source{petersen:page-0352}
  \uses{thm:pet-ch9-bochner-vanishing-1forms, def:pet-ch9-hodge-cohomology}
  \dcref{ch9:9.2.5}
  For a compact $n$-manifold with nonnegative Ricci curvature,
  \[
    b_1(M)\le n.
  \]
  Equality holds precisely for flat tori.
\end{corollary}
\begin{proof}
  \uses{cor:pet-ch9-betti1-bound-flat-torus}
  \sketch
  Hodge theory and \cref{thm:pet-ch9-bochner-vanishing-1forms} identify
  \(H^1(M)\) with a space of parallel forms.  Evaluation at one point is
  injective, so \(b_1\le n\).  Equality supplies a global parallel frame and
  forces flatness.  On the universal cover \(\mathbb R^n\), deck
  transformations preserve that frame and are therefore translations;
  compactness and discreteness make the translation group a full lattice
  \(\mathbb Z^n\).  Conversely, a flat torus has \(n\) independent parallel
  one-forms.
\end{proof}

Thus \cref{cor:pet-ch9-betti1-bound-flat-torus} obstructs nonnegative Ricci
curvature whenever a closed $n$-manifold has first Betti number greater than
$n$.

\subsection{The Estimation Theorem}

Under a negative Ricci lower bound, the Bochner inequality combines with a
Sobolev estimate to control the dimension of the harmonic-form space.

\begin{lemma}[Li's dimension estimate]
  \label{lem:pet-ch9-li-dimension-estimate}
  \lean{PetersenLib.li_dimension_estimate}
  \source{petersen:page-0353,petersen:page-0354}
  \dcref{ch9:9.2.6}
  Let $E\to M$ be a rank-$m$ metric vector bundle over compact $(M,g)$.  For
  $s\in\Gamma(E)$ define
  \[
    \|s\|_\infty=\max_M|s|,
    \qquad
    \|s\|_p=\left(\frac1{\operatorname{vol}M}
      \int_M|s|^p\,\mathrm{vol}\right)^{1/p}.
  \]
  If $V\subset\Gamma(E)$ is finite dimensional and
  \[
    C(V)=\max_{0\ne s\in V}\frac{\|s\|_\infty}{\|s\|_2},
  \]
  then $\dim V\le mC(V)^2$.
\end{lemma}
\begin{proof}
  \uses{lem:pet-ch9-li-dimension-estimate}
  \sketch
  For an \(L^2\)-orthonormal basis \(e_1,\ldots,e_\ell\) of \(V\), put
  \(F(x)=\sum_i|e_i(x)|^2\).  Its normalized integral is \(\ell\).  At a
  maximum \(x_0\), choose the basis so that only the first
  \(r=\operatorname{rank}(V\to E_{x_0})\le m\) elements are nonzero.  The
  defining supremum estimate gives \(|e_i(x_0)|^2\le C(V)^2\) for each
  nonzero term, and hence
  \(\ell\le F(x_0)\le mC(V)^2\).
\end{proof}

\begin{theorem}[Moser iteration supremum estimate]
  \label{thm:pet-ch9-moser-iteration}
  \lean{PetersenLib.moserIteration}
  \source{petersen:page-0354,petersen:page-0355}
  \dcref{ch9:9.2.7}
  Assume compact $(M,g)$ satisfies, for some $\nu>1$ and every smooth $u$,
  \[
    \|u\|_{2\nu}\le S\|\nabla u\|_2+\|u\|_2.
  \]
  If a continuous $f:M\to[0,\infty)$ is smooth where positive and obeys
  $\Delta f\ge-\lambda f$, then
  \[
    \|f\|_\infty\le\exp\Big(\frac{S\sqrt{\lambda\nu}}{\sqrt{\nu-1}}\Big)\|f\|_2.
  \]
\end{theorem}
\begin{proof}
  \uses{thm:pet-ch9-moser-iteration}
  \sketch
  The differential inequality holds distributionally across \(\{f=0\}\).
  Testing it against \(f^{2q-1}\), \(q\ge1\), gives
  \[
    \|d(f^q)\|_2^2\le\frac{q^2\lambda}{2q-1}\|f^q\|_2^2.
  \]
  Applying the Sobolev inequality to \(f^q\) yields
  \[
    \|f\|_{2q\nu}
    \le\left(1+Sq\sqrt{\frac{\lambda}{2q-1}}\right)^{1/q}\|f\|_{2q}.
  \]
  Iterate with \(q=\nu^k\).  Taking logarithms, using
  \(\log(1+x)\le x\), and summing the resulting geometric series bounds the
  product by
  \(\exp(S\sqrt{\lambda\nu}/\sqrt{\nu-1})\), as claimed.
\end{proof}

\begin{lemma}[Kato's inequality]
  \label{lem:pet-ch9-kato-inequality}
  \lean{PetersenLib.katoInequality}
  \source{petersen:page-0356}
  Let $f=|T|$ for a tensor field $T$.  At every point where $f$ is smooth,
  \[
    df\le|\nabla T|.
  \]
\end{lemma}
\begin{proof}
  \uses{lem:pet-ch9-kato-inequality}
  \sketch
  Where \(|T|>0\), differentiating \(|T|^2\) and applying Cauchy--Schwarz gives
  \(2|T|\,d|T|=2g(\nabla T,T)\le2|\nabla T|\,|T|\).  Division by
  \(2|T|\) proves the estimate.
\end{proof}

\begin{theorem}[Gromov--Gallot first Betti estimate]
  \label{thm:pet-ch9-gromov-gallot-betti1-estimate}
  \lean{PetersenLib.gromovGallot_firstBettiNumber_estimate}
  \source{petersen:page-0355,petersen:page-0356}
  \uses{lem:pet-ch9-li-dimension-estimate, thm:pet-ch9-moser-iteration, lem:pet-ch9-kato-inequality, def:pet-ch9-hodge-cohomology, prop:pet-ch9-bochner-formula-1forms}
  \dcref{ch9:9.2.8}
  Let $M^n$ be compact with
  $\operatorname{Ric}\ge(n-1)k$ and $\operatorname{diam}M\le D$.  There is a
  bound depending only on $n$ and $kD^2$,
  \[
    b_1(M)\le C(n,kD^2),
    \qquad \lim_{\varepsilon\to0}C(n,\varepsilon)=n.
  \]
  Hence some $\varepsilon(n)>0$ satisfies
  $kD^2\ge-\varepsilon(n)\Longrightarrow b_1(M)\le n$.
\end{theorem}
\begin{proof}
  \uses{thm:pet-ch9-gromov-gallot-betti1-estimate}
  \sketch
  For a harmonic one-form \(\omega\), set \(f=|\omega|\).  Kato's inequality
  from \cref{lem:pet-ch9-kato-inequality} and the identity in
  \cref{prop:pet-ch9-bochner-formula-1forms} imply
  \(\Delta f\ge(n-1)kf\).  The Sobolev constant controlled by \(n,kD^2\),
  together with \cref{thm:pet-ch9-moser-iteration}, bounds
  \(\|\omega\|_\infty/\|\omega\|_2\) by a function tending to \(1\) as
  \(kD^2\to0\).  Applying \cref{lem:pet-ch9-li-dimension-estimate} to
  \(\mathcal H^1(M)\) gives the asserted bound and its limiting value \(n\).
\end{proof}

\section{Lichnerowicz Laplacians}

The connection and Lichnerowicz Laplacians extend the preceding argument to
tensor bundles.

\subsection{The Connection Laplacian}

Fix a tensor bundle $E\to M$ of fibre dimension $m$ (e.g.\ $p$-forms,
symmetric tensors, or curvature tensors).

\begin{proposition}[maximum principle for the connection Laplacian]
  \label{prop:pet-ch9-vanishing-connection-laplacian}
  \lean{PetersenLib.vanishing_connectionLaplacian}
  \source{petersen:page-0357}
  \dcref{ch9:9.3.1}
  Suppose a section $T\in\Gamma(E)$ over a Riemannian manifold satisfies
  $g(\nabla^*\nabla T,T)\le0$ pointwise.  If the function $|T|$ has a global
  maximum, then $\nabla T=0$.
\end{proposition}
\begin{proof}
  \uses{prop:pet-ch9-vanishing-connection-laplacian}
  \sketch
  The tensor Bochner identity gives
  \(\Delta\tfrac12|T|^2=|\nabla T|^2-g(\nabla^*\nabla T,T)\ge0\).
  At a maximum of \(|T|\), the maximum principle forces this function to be
  constant and hence \(\nabla T=0\).
\end{proof}

\begin{theorem}[Gallot's connection-Laplacian estimate]
  \label{thm:pet-ch9-gallot-connection-laplacian-estimate}
  \lean{PetersenLib.gallot_connectionLaplacian_estimate}
  \source{petersen:page-0357}
  \uses{lem:pet-ch9-li-dimension-estimate, thm:pet-ch9-moser-iteration, lem:pet-ch9-kato-inequality}
  \dcref{ch9:9.3.2}
  Let compact $(M,g)$ satisfy the Sobolev hypothesis in
  \cref{thm:pet-ch9-moser-iteration}, and let $V\subset\Gamma(E)$ be finite
  dimensional.  If every $T\in V$ obeys
  $g(\nabla^*\nabla T,T)\le\lambda|T|^2$, then
  \[
    \dim V\le
    m\cdot\exp\Big(\frac{2S\sqrt{\lambda\nu}}{\sqrt{\nu-1}}\Big).
  \]
\end{theorem}
\begin{proof}
  \uses{thm:pet-ch9-gallot-connection-laplacian-estimate}
  \sketch
  For \(f=|T|\), the tensor Bochner identity and
  \cref{lem:pet-ch9-kato-inequality} give \(\Delta f\ge-\lambda f\).
  The argument of \cref{thm:pet-ch9-gromov-gallot-betti1-estimate}, now using
  \cref{thm:pet-ch9-moser-iteration}, bounds
  \(C(V)\) by
  \(\exp(S\sqrt{\lambda\nu}/\sqrt{\nu-1})\). Squaring this bound in
  \cref{lem:pet-ch9-li-dimension-estimate} gives the conclusion.
\end{proof}

\subsection{The Weitzenböck Curvature}

\begin{definition}[Weitzenböck curvature operator]
  \label{def:pet-ch9-weitzenbock-curvature-operator}
  \lean{PetersenLib.weitzenbockCurvatureOperator}
  \source{petersen:page-0357,petersen:page-0358}
  Given a $(0,k)$-tensor $T$, define its \emph{Weitzenb\"ock curvature} by
  \[
    \operatorname{Ric}(T)(X_1,\dots,X_k)=\sum_j\big(R(e_j,X_i)T\big)(X_1,\dots,e_j,\dots,X_k)
  \]
  where $e_j$ is orthonormal and the displayed slot entails summation over
  $i$.  For $\omega=X^\flat$ this specializes to
  \[
    \operatorname{Ric}(\omega)(X')
      =\sum_j(R(e_j,X')\omega)(e_j)
      =-\omega\left(\sum_jR(e_j,X')e_j\right)
      =\omega(\operatorname{Ric}(X')).
  \]
  Thus the notation agrees with the ordinary Ricci endomorphism on one-forms
  (the same operator is also denoted $W$).
\end{definition}

\begin{definition}[Lichnerowicz Laplacian]
  \label{def:pet-ch9-lichnerowicz-laplacian}
  \lean{PetersenLib.lichnerowiczLaplacian}
  \source{petersen:page-0358}
  \uses{def:pet-ch9-weitzenbock-curvature-operator}
  For $c>0$, the \emph{Lichnerowicz Laplacian} on $E$ is
  \[
    \Delta_L=\nabla^*\nabla+c\operatorname{Ric}.
  \]
  Its Hodge-form instance has $c=1$; for symmetric $(0,2)$-tensors and for
  the curvature tensor, take $c=\tfrac12$.
\end{definition}

The Bochner technique applies to $T$ with $\Delta_LT=0$: the maximum principle
(\cref{prop:pet-ch9-vanishing-connection-laplacian}) forces $T$ parallel once
$g(\nabla^*\nabla T,T)\le0$, which by $\Delta_LT=0$ is equivalent to
$g(\operatorname{Ric}(T),T)\ge0$. The remaining work is to reformulate
$g(\operatorname{Ric}(T),T)$ so that nonnegativity of the curvature operator makes this
manifest, and to identify natural Lichnerowicz Laplacians.

\subsection{Simplification of $\operatorname{Ric}(T)$}

\begin{convention}[metric on $\mathfrak{so}(T_pM)$]
  \label{conv:pet-ch9-so-metric-convention}
  \lean{PetersenLib.soMetricConvention}
  \source{petersen:page-0359}
  \uses{def:pet-ch9-weitzenbock-curvature-operator}
  Identify $\Lambda^2T_pM$ isometrically with $\mathfrak{so}(T_pM)$ by sending
  a unit bivector $v\wedge w$ to the quarter-turn in
  $\operatorname{span}\{v,w\}$.  This uses the transported bivector metric,
  rather than the Euclidean operator norm for which that quarter-turn has
  length $\sqrt2$.  If $\Xi_\sigma$ is orthonormal for this metric and
  $\mathfrak R$ is the curvature operator, then the skew map
  $R_{X,Y}$ expands as
  \[
    R_{X,Y}=\sum_\sigma g(R_{X,Y},\Xi_\sigma)\Xi_\sigma
    =\sum_\sigma g(\mathfrak{R}(X\wedge Y),\Xi_\sigma)\Xi_\sigma
    =\sum_\sigma g(\mathfrak{R}(\Xi_\sigma),X\wedge Y)\Xi_\sigma
    =-\sum_\sigma g(R(\Xi_\sigma)X,Y)\Xi_\sigma.
  \]
\end{convention}

\begin{lemma}[curvature-operator formula for Weitzenb\"ock curvature]
  \label{lem:pet-ch9-ric-curvature-operator-formula}
  \lean{PetersenLib.weitzenbock_eq_curvatureOperator_sum}
  \source{petersen:page-0359,petersen:page-0360}
  \uses{conv:pet-ch9-so-metric-convention, def:pet-ch9-weitzenbock-curvature-operator, def:pet-ch9-lichnerowicz-laplacian}
  \dcref{ch9:9.3.3}
  Every $(0,k)$-tensor $T$ satisfies
  \[
    \operatorname{Ric}(T)=-\sum_\sigma R(\Xi_\sigma)(\Xi_\sigma T),\qquad
    \Delta_LT=\nabla^*\nabla T-c\sum_\sigma R(\Xi_\sigma)(\Xi_\sigma T).
  \]
  The endomorphism $T\mapsto\operatorname{Ric}(T)$ is self-adjoint as well.
\end{lemma}
\begin{proof}
  \uses{lem:pet-ch9-ric-curvature-operator-formula}
  \sketch
  Expand each \(R(e_j,X_i)\) in the orthonormal basis
  \(\Xi_\sigma\) using \cref{conv:pet-ch9-so-metric-convention}.  Resumming
  the \(e_j\)-components in the \(i\)-th slot yields
  \(\operatorname{Ric}(T)=-\sum_\sigma R(\Xi_\sigma)(\Xi_\sigma T)\), and
  substitution gives the formula for \(\Delta_L\).  If
  \(\mathfrak R(\Xi_\sigma)=\lambda_\sigma\Xi_\sigma\), skew-adjointness of
  \(\Xi_\sigma\) gives
  \[
    g(\operatorname{Ric}(T),S)
      =\sum_\sigma\lambda_\sigma g(\Xi_\sigma T,\Xi_\sigma S),
  \]
  which is symmetric in \(T,S\).
\end{proof}

\begin{corollary}[Weitzenb\"ock lower bound]
  \label{cor:pet-ch9-ric-nonneg-curvature-operator}
  \lean{PetersenLib.weitzenbock_nonneg_of_curvatureOperator_nonneg}
  \source{petersen:page-0360}
  \uses{lem:pet-ch9-ric-curvature-operator-formula}
  \dcref{ch9:9.3.4}
  Nonnegative curvature operator implies
  $g(\operatorname{Ric}(T),T)\ge0$ for every tensor $T$.  More generally, from
  $\mathfrak R\ge k$ with $k<0$ one obtains
  \[
    g(\operatorname{Ric}(T),T)\ge kC|T|^2,
  \]
  where $C>0$ depends only on the tensor type.
\end{corollary}
\begin{proof}
  \uses{cor:pet-ch9-ric-nonneg-curvature-operator}
  \sketch
  Diagonalizing \(\mathfrak R\) as in
  \cref{lem:pet-ch9-ric-curvature-operator-formula} gives
  \[
    g(\operatorname{Ric}(T),T)=\sum_\sigma\lambda_\sigma|\Xi_\sigma T|^2.
  \]
  This is nonnegative when every \(\lambda_\sigma\ge0\).  If
  \(\lambda_\sigma\ge k<0\), the tensor-type estimate
  \(\sum_\sigma|\Xi_\sigma T|^2\le C|T|^2\) gives the stated lower bound.
\end{proof}

\begin{theorem}[dimension estimate for Lichnerowicz-harmonic tensors]
  \label{thm:pet-ch9-lichnerowicz-dimension-estimate}
  \lean{PetersenLib.lichnerowiczLaplacian_kernel_dimension_estimate}
  \source{petersen:page-0360}
  \uses{cor:pet-ch9-ric-nonneg-curvature-operator, thm:pet-ch9-gallot-connection-laplacian-estimate, prop:pet-ch9-vanishing-connection-laplacian}
  \dcref{ch9:9.3.5}
  Suppose $\mathfrak R\ge k$ and $\operatorname{diam}M\le D$, and set
  \[
    V=\ker\Delta_L
      =\{T\in\Gamma(E):\nabla^*\nabla T+c\operatorname{Ric}(T)=0\}.
  \]
  Then
  \[
    \dim V\le
    m\cdot\exp\Big(2D\cdot C(n,kD^2)
      \frac{\sqrt{-kcC\nu}}{\sqrt{\nu-1}}\Big).
  \]
  If $k=0$, all elements of $V$ are parallel.
\end{theorem}

\section{The Bochner Technique in General}

\subsection{Forms}

\begin{theorem}[Weitzenb\"ock formula for the Hodge Laplacian]
  \label{thm:pet-ch9-weitzenbock-hodge-laplacian}
  \lean{PetersenLib.hodgeLaplacian_eq_lichnerowiczLaplacian}
  \source{petersen:page-0361}
  \uses{def:pet-ch9-hodge-laplacian, def:pet-ch9-weitzenbock-curvature-operator, def:pet-ch9-lichnerowicz-laplacian}
  \dcref{ch9:9.4.1}
  On differential forms, the Lichnerowicz operator with coefficient one is
  exactly the Hodge Laplacian:
  \[
    \square\omega=(d\delta+\delta d)\omega=\nabla^*\nabla\omega+\operatorname{Ric}(\omega).
  \]
\end{theorem}
\begin{proof}
  \uses{thm:pet-ch9-weitzenbock-hodge-laplacian}
  \sketch
  In a normal orthonormal frame, expand
  \(\delta\omega=-\sum_i i_{E_i}\nabla_{E_i}\omega\) and the alternating
  formula for \(d\omega\).  In \(d\delta\omega+\delta d\omega\), the equal-order
  terms combine to \(\nabla^*\nabla\omega\); commuting the remaining second
  derivatives produces
  \(\sum_{i,j}(R(E_j,X_i)\omega)(X_1,\ldots,E_j,\ldots,X_k)\), which is
  \(\operatorname{Ric}(\omega)\).
\end{proof}

\subsection{The Curvature Tensor}

\begin{theorem}[Lichnerowicz identity for the curvature tensor]
  \label{thm:pet-ch9-curvature-tensor-lichnerowicz}
  \lean{PetersenLib.curvatureTensor_lichnerowiczIdentity}
  \source{petersen:page-0362,petersen:page-0363}
  \uses{def:pet-ch9-lichnerowicz-laplacian, def:pet-ch9-weitzenbock-curvature-operator}
  \dcref{ch9:9.4.2}
  For the curvature tensor, the Lichnerowicz operator with coefficient
  $\tfrac12$ satisfies
  \begin{align*}
    (\nabla^*\nabla R)(X,Y,Z,W)+\tfrac12\operatorname{Ric}(R)(X,Y,Z,W)
    &=\tfrac12(\nabla_X\nabla^*R)(Y,Z,W)-\tfrac12(\nabla_Y\nabla^*R)(X,Z,W)\\
    &\quad+\tfrac12(\nabla_Z\nabla^*R)(W,X,Y)-\tfrac12(\nabla_W\nabla^*R)(Z,X,Y).
  \end{align*}
  Thus the $c=\tfrac12$ Lichnerowicz Laplacian of $R$ vanishes whenever
  $\nabla^*R=0$; curvature is Lichnerowicz-harmonic exactly in this
  divergence-free case.
\end{theorem}
\begin{proof}
  \uses{thm:pet-ch9-curvature-tensor-lichnerowicz}
  \sketch
  Evaluate in a normal frame and apply the second Bianchi identity to the two
  differentiated curvature slots.  Commuting the resulting second covariant
  derivatives contributes the Weitzenb\"ock term.  Repeating the computation
  after interchanging the pair \((X,Y)\) with \((Z,W)\), and averaging by the
  pair symmetry of \(R\), gives exactly the four divergence derivatives on
  the right and the coefficient \(\tfrac12\operatorname{Ric}(R)\).  If
  \(\nabla^*R=0\), all four terms vanish.
\end{proof}

\subsection{Symmetric $(0,2)$-Tensors}

\begin{definition}[the exterior derivative $d^\nabla h$ of a symmetric tensor]
  \label{def:pet-ch9-codazzi-exterior-derivative}
  \lean{PetersenLib.symmetricTensorExteriorDerivative}
  \source{petersen:page-0363}
  If $h(z,w)=g(H(z),w)$ is symmetric, then the connection exterior derivative
  is defined by
  \[
    d^\nabla h(X,Y,Z)=(\nabla_Xh)(Y,Z)-(\nabla_Yh)(X,Z).
  \]
\end{definition}

\begin{remark}[Codazzi identity for the Ricci tensor]
  \label{rem:pet-ch9-ricci-codazzi-identity}
  \lean{PetersenLib.ricciTensor_codazziIdentity}
  \source{petersen:page-0363}
  \uses{def:pet-ch9-codazzi-exterior-derivative}
  The second fundamental form of an immersed hypersurface
  $M^n\to\mathbb R^{n+1}$ satisfies the Codazzi--Mainardi equation
  $d^\nabla\mathrm{II}=0$. For the Ricci tensor,
  \[
    (d^\nabla\operatorname{Ric})(X,Y,Z)=(\nabla^*R)(Z,X,Y),
  \]
  and the Schouten--Weyl analogue is the corresponding identity from
  Petersen, exercise 3.4.26.
\end{remark}

\begin{theorem}[Lichnerowicz identity for symmetric tensors]
  \label{thm:pet-ch9-symmetric-tensor-lichnerowicz-identity}
  \lean{PetersenLib.symmetricTensor_lichnerowiczIdentity}
  \source{petersen:page-0364}
  \uses{def:pet-ch9-codazzi-exterior-derivative, def:pet-ch9-weitzenbock-curvature-operator}
  \dcref{ch9:9.4.3}
  For every symmetric $(0,2)$-tensor $h$ and vector $X$,
  \[
    (\nabla_X\nabla^*h)(X)+(\nabla^*d^\nabla h)(X,X)
    =(\nabla^*\nabla h)(X,X)+\tfrac12(\operatorname{Ric}(h))(X,X).
  \]
\end{theorem}
\begin{proof}
  \uses{thm:pet-ch9-symmetric-tensor-lichnerowicz-identity}
  \sketch
  At the evaluation point use a normal frame to expand
  \(\nabla_X\nabla^*h\) and \(\nabla^*d^\nabla h\).  Their second-derivative
  terms combine to \(\nabla^*\nabla h\); commuting the remaining derivatives
  gives \((R(E_i,X)h)(E_i,X)\).  Symmetry of \(h\) identifies this contraction
  with \(\tfrac12\operatorname{Ric}(h)(X,X)\).
\end{proof}

\begin{definition}[Codazzi and harmonic symmetric tensors]
  \label{def:pet-ch9-codazzi-harmonic-tensor}
  \lean{PetersenLib.IsCodazziTensor,PetersenLib.IsHarmonicSymmetricTensor}
  \source{petersen:page-0364}
  \uses{def:pet-ch9-codazzi-exterior-derivative}
  Call a symmetric tensor $h$ \emph{Codazzi} when $d^\nabla h=0$; call it
  \emph{harmonic} when it is Codazzi and also divergence free,
  $\nabla^*h=0$.
\end{definition}

\begin{proposition}[harmonic tensors and constant trace]
  \label{prop:pet-ch9-harmonic-iff-codazzi-constant-trace}
  \lean{PetersenLib.harmonic_iff_codazzi_constantTrace}
  \source{petersen:page-0364}
  \uses{def:pet-ch9-codazzi-harmonic-tensor}
  \dcref{ch9:9.4.4}
  For a symmetric $(0,2)$-tensor $h$,
  \[
    h\text{ is harmonic}\quad\Longleftrightarrow\quad
    d^\nabla h=0\text{ and }\operatorname{tr}h\text{ is constant}.
  \]
\end{proposition}
\begin{proof}
  \uses{prop:pet-ch9-harmonic-iff-codazzi-constant-trace}
  \sketch
  In a normal frame, symmetry of \(h\) gives
  \[
    (\nabla^*h)(X)=-D_X\operatorname{tr}h
      +\sum_i(d^\nabla h)(X,E_i,E_i).
  \]
  For a Codazzi tensor the final sum vanishes, so divergence freedom is
  equivalent to constancy of the trace.
\end{proof}

In particular, a constant-mean-curvature hypersurface has harmonic second
fundamental form.  By \cref{rem:pet-ch9-ricci-codazzi-identity}, a Codazzi Ricci
tensor is equivalent to divergence-free curvature.

\begin{corollary}[harmonic curvature and Codazzi Ricci tensor]
  \label{cor:pet-ch9-ricci-harmonic-iff-curvature-harmonic}
  \lean{PetersenLib.ricciTensor_harmonic_iff_curvatureTensor_harmonic}
  \source{petersen:page-0365}
  \uses{prop:pet-ch9-harmonic-iff-codazzi-constant-trace, rem:pet-ch9-ricci-codazzi-identity, thm:pet-ch9-curvature-tensor-lichnerowicz}
  \dcref{ch9:9.4.5}
  The following conditions are equivalent: $\operatorname{Ric}$ is harmonic,
  $\operatorname{Ric}$ is Codazzi with constant trace, and
  $\nabla^*R=0$ (equivalently, $R$ is harmonic).
\end{corollary}
\begin{proof}
  \uses{cor:pet-ch9-ricci-harmonic-iff-curvature-harmonic}
  \sketch
  The identity in \cref{rem:pet-ch9-ricci-codazzi-identity} equates
  \(d^\nabla\operatorname{Ric}\) with the divergence of \(R\), while the
  contracted Bianchi identity makes \(\nabla^*\operatorname{Ric}\) a constant
  multiple of \(d\operatorname{scal}\).  Thus \(\operatorname{Ric}\) is
  harmonic exactly when \(R\) is divergence free.  The constant-trace
  characterization in \cref{prop:pet-ch9-harmonic-iff-codazzi-constant-trace}
  yields the equivalent formulations.
\end{proof}

\subsection{Topological and Geometric Consequences}

\begin{theorem}[forms under a curvature-operator lower bound]
  \label{thm:pet-ch9-nonneg-curvature-operator-parallel-forms}
  \lean{PetersenLib.nonnegCurvatureOperator_harmonicForms_parallel}
  \source{petersen:page-0365}
  \uses{thm:pet-ch9-weitzenbock-hodge-laplacian, thm:pet-ch9-lichnerowicz-dimension-estimate, cor:pet-ch9-ric-nonneg-curvature-operator}
  \dcref{ch9:9.4.6}
  Let $(M,g)$ be closed of dimension $n$.  Under $\mathfrak R\ge0$, every
  harmonic form is parallel.  If $\mathfrak R>0$, a parallel $l$-form can be
  nonzero only for $l=0$ or $l=n$.  For the lower bound $\mathfrak R\ge k$ and
  diameter bound $D$,
  \[
    b_l(M)\le\binom{n}{l}
      \exp\big(2D\cdot C(n,kD^2)\sqrt{-kC}\big).
  \]
\end{theorem}
\begin{proof}
  \uses{thm:pet-ch9-nonneg-curvature-operator-parallel-forms}
  \sketch
  By \cref{thm:pet-ch9-weitzenbock-hodge-laplacian} and
  \cref{cor:pet-ch9-ric-nonneg-curvature-operator}, a harmonic form satisfies
  \(g(\nabla^*\nabla\omega,\omega)\le0\); compactness and
  \cref{prop:pet-ch9-vanishing-connection-laplacian} make it parallel.  If
  \(\mathfrak R>0\), the eigenbasis formula forces
  \(L\omega=0\) for every skew endomorphism \(L\), which annihilates all
  degrees except \(0\) and \(n\).  Finally apply
  \cref{thm:pet-ch9-lichnerowicz-dimension-estimate} to the Hodge Laplacian on
  \(l\)-forms, whose fibre rank is \(\binom nl\).
\end{proof}

\begin{example}[nonnegative sectional curvature on $\mathbb{CP}^2\#\mathbb{CP}^2$]
  \label{ex:pet-ch9-cp2-connect-sum-nonneg-curvature}
  \lean{PetersenLib.cp2ConnectSumCp2_nonnegSectionalCurvature}
  \source{petersen:page-0366}
  \uses{thm:pet-ch9-nonneg-curvature-operator-parallel-forms}
  \dcref{ch9:9.4.7}
  Let $S^1$ act diagonally on $S^2\times S^3$, by rotation of total angle
  $2\pi k$ on $S^2$ and by the Hopf action on
  $S^3\subset\mathbb C^2$.  The quotient is $S^2\times S^2$ for even $k$ and
  $\mathbb{CP}^2\#\mathbb{CP}^2$ for odd $k$; O'Neill's formula gives it
  nonnegative sectional curvature.  Petersen's theorem 10.3.7 restricts the
  simply connected topological $4$-manifolds with $\mathfrak R\ge0$ to
  $S^2\times S^2$, $S^4$, and $\mathbb{CP}^2$.  Hence the odd quotient
  distinguishes $\sec\ge0$ from $\mathfrak R\ge0$.
\end{example}

\begin{theorem}[Tachibana rigidity]
  \label{thm:pet-ch9-tachibana-parallel-curvature}
  \lean{PetersenLib.tachibana_parallelCurvature}
  \source{petersen:page-0367}
  \uses{thm:pet-ch9-curvature-tensor-lichnerowicz, prop:pet-ch9-vanishing-connection-laplacian, conv:pet-ch9-so-metric-convention}
  \dcref{ch9:9.4.8}
  On a closed manifold, divergence-free curvature together with
  $\mathfrak R\ge0$ implies $\nabla R=0$.  If the curvature operator is
  positive, the metric has constant sectional curvature.
\end{theorem}
\begin{proof}
  \uses{thm:pet-ch9-tachibana-parallel-curvature}
  \sketch
  The curvature identity \cref{thm:pet-ch9-curvature-tensor-lichnerowicz}, the
  lower bound in \cref{cor:pet-ch9-ric-nonneg-curvature-operator}, and
  \cref{prop:pet-ch9-vanishing-connection-laplacian} imply \(\nabla R=0\).
  Under \(\mathfrak R>0\), the argument from
  \cref{thm:pet-ch9-nonneg-curvature-operator-parallel-forms} gives
  \(LR=0\) for every \(L\in\mathfrak{so}(T_pM)\).  Choosing rotations that
  exchange an arbitrary vector with an orthogonal one and polarizing shows
  all mixed curvature components vanish.  Every decomposable bivector is
  consequently an eigenvector of \(\mathfrak R\), forcing
  \(\mathfrak R=kI\).
\end{proof}

\subsection{Simplification of $g(\operatorname{Ric}(T),T)$}

\begin{definition}[$\Lambda^2TM$-valued tensor $\hat T$]
  \label{def:pet-ch9-tensor-lambda2-valued}
  \lean{PetersenLib.lambda2ValuedTensor}
  \source{petersen:page-0368}
  \uses{conv:pet-ch9-so-metric-convention}
  Associate to a $(0,k)$-tensor $T$ the
  $\Lambda^2TM\cong\mathfrak{so}(TM)$-valued tensor $\hat T$ uniquely
  characterized by
  \[
    g\big(L,\hat T(X_1,\dots,X_k)\big)=(LT)(X_1,\dots,X_k)\quad\text{for all }
    L\in\mathfrak{so}(TM).
  \]
\end{definition}

\begin{lemma}[hat-tensor formula for Weitzenb\"ock curvature]
  \label{lem:pet-ch9-ric-curvature-operator-hat-formula}
  \lean{PetersenLib.weitzenbock_eq_curvatureOperator_hat}
  \source{petersen:page-0368}
  \uses{def:pet-ch9-tensor-lambda2-valued, def:pet-ch9-weitzenbock-curvature-operator, lem:pet-ch9-ric-curvature-operator-formula}
  \dcref{ch9:9.4.9}
  For any two $(0,k)$-tensors $T$ and $S$, contraction of their hat tensors
  gives
  \[
    g(\operatorname{Ric}(T),S)=g(\mathfrak R(\hat T),\hat S).
  \]
\end{lemma}
\begin{proof}
  \uses{lem:pet-ch9-ric-curvature-operator-hat-formula}
  \sketch
  In orthonormal bases \(e_I\) and \(\Xi_\sigma\), use
  \cref{lem:pet-ch9-ric-curvature-operator-formula} and the definition in
  \cref{def:pet-ch9-tensor-lambda2-valued} to write
  \[
    g(\operatorname{Ric}(T),S)
      =\sum_{\sigma,I}g(\Xi_\sigma,\hat T(e_I))
        g(R(\Xi_\sigma),\hat S(e_I)).
  \]
  Resumming the \(\Xi_\sigma\)-components is precisely
  \(g(\mathfrak R(\hat T),\hat S)\).
\end{proof}

\begin{proposition}[hat tensor of a $1$-form]
  \label{prop:pet-ch9-1form-hat-ricci}
  \lean{PetersenLib.oneForm_hat_ricci_identity}
  \source{petersen:page-0368,petersen:page-0369}
  \uses{lem:pet-ch9-ric-curvature-operator-hat-formula, def:pet-ch9-tensor-lambda2-valued}
  \dcref{ch9:9.4.10}
  For a one-form $\omega=X^\flat$,
  \[
    \hat\omega(Z)=X\wedge Z,
    \qquad
    g(\mathfrak R(\hat\omega),\hat\omega)=\operatorname{Ric}(X,X).
  \]
\end{proposition}
\begin{proof}
  \uses{prop:pet-ch9-1form-hat-ricci}
  \sketch
  For skew \(L\),
  \((L\omega)(Z)=-g(X,LZ)=g(L,X\wedge Z)\), hence
  \(\hat\omega(Z)=X\wedge Z\).  Applying
  \cref{lem:pet-ch9-ric-curvature-operator-hat-formula} and tracing over an
  orthonormal frame gives
  \(g(\mathfrak R(\hat\omega),\hat\omega)=
    \sum_iR(X,E_i,E_i,X)=\operatorname{Ric}(X,X)\).
\end{proof}

\begin{remark}[$p$-form generalization]
  \label{rem:pet-ch9-pform-hat-formula}
  \lean{PetersenLib.pForm_hat_formula}
  \source{petersen:page-0369}
  \uses{def:pet-ch9-tensor-lambda2-valued, prop:pet-ch9-1form-hat-ricci}
  For a $p$-form $\omega$, define the vector-valued $(p-1)$-form $\Omega$ by
  \(g(\Omega(X_1,\dots,X_{p-1}),X_p)=\omega(X_1,\dots,X_p)\).  Then
  \[
    \hat\omega(X_1,\dots,X_p)=\sum_{i=1}^p(-1)^{p-i}X_i\wedge
    \Omega(X_1,\dots,\hat X_i,\dots,X_p),
  \]
  Moreover, $\hat\omega=0$ at a point only if $\omega=0$ there.
\end{remark}

Turning to $(0,2)$-tensors, let $h(z,w)=g(H(z),w)$ with adjoint $H^*$.

\begin{proposition}[hat tensor of a $(0,2)$-tensor]
  \label{prop:pet-ch9-02tensor-hat-formula}
  \lean{PetersenLib.symmetricTensor_hat_formula}
  \source{petersen:page-0369,petersen:page-0370}
  \uses{def:pet-ch9-tensor-lambda2-valued}
  \dcref{ch9:9.4.11}
  If $h(z,w)=g(H(z),w)$ and $H^*$ is the adjoint of $H$, then
  \[
    \hat h(z,w)=H(z)\wedge w-z\wedge H^*(w).
  \]
  The tensor $\hat h$ vanishes exactly when $h$ is a scalar multiple of $g$.
\end{proposition}
\begin{proof}
  \uses{prop:pet-ch9-02tensor-hat-formula}
  \sketch
  For skew \(L\), expanding \((Lh)(z,w)\) and moving \(L\) through the metric
  yields
  \((Lh)(z,w)=g(L,H(z)\wedge w-z\wedge H^*(w))\), proving the formula.
  Scalar multiples of \(g\) clearly have zero hat tensor.  Conversely,
  \(\hat h=0\) implies
  \(H(z)\wedge w=z\wedge H^*(w)\) for all \(z,w\); varying the pair forces
  every vector to be an eigenvector of \(H\) with one common eigenvalue, so
  \(H=\lambda I\).
\end{proof}

If $H$ is normal it diagonalizes over $T_pM\otimes\mathbb{C}$: say $H(z)=\lambda z$,
$H(w)=\mu w$ for orthonormal $z,w\in T_pM\otimes\mathbb{C}$. Then
\[
  g\big(\mathfrak{R}(H(z)\wedge w-z\wedge H^*(w)),H(z)\wedge w-z\wedge H^*(w)\big)
  =|\lambda-\bar\mu|^2\,g\big(\mathfrak{R}(\bar z\wedge w),\bar z\wedge w\big).
\]

\begin{definition}[complex sectional curvature]
  \label{def:pet-ch9-complex-sectional-curvature}
  \lean{PetersenLib.complexSectionalCurvature}
  \source{petersen:page-0370}
  \uses{prop:pet-ch9-02tensor-hat-formula}
  For complexified tangent vectors $z,w\in T_pM\otimes\mathbb C$, define the
  \emph{complex sectional curvature} of their span to be
  \[
    g(\mathfrak R(z\wedge w),\bar z\wedge\bar w).
  \]
\end{definition}

\begin{lemma}[real expansion of the complex sectional curvature]
  \label{lem:pet-ch9-complex-sectional-curvature-real-formula}
  \lean{PetersenLib.complexSectionalCurvature_realFormula}
  \source{petersen:page-0370,petersen:page-0371}
  \uses{def:pet-ch9-complex-sectional-curvature}
  If $z=x+iy$ and $w=u+iv$, with $x,y,u,v\in T_pM$, then
  \begin{align*}
    g(\mathfrak{R}(z\wedge w),\bar z\wedge\bar w)
    &=R(x,u,u,x)+R(y,v,v,y)+R(x,v,v,x)+R(y,u,u,y)+2R(y,x,v,u).
  \end{align*}
  Thus $\mathfrak R\ge0$ implies nonnegative complex sectional curvature,
  which in turn implies nonnegative real sectional curvature; neither
  converse is asserted.
\end{lemma}
\begin{proof}
  \uses{lem:pet-ch9-complex-sectional-curvature-real-formula}
  \sketch
  Write
  \(z\wedge w=(x\wedge u-y\wedge v)+i(x\wedge v+y\wedge u)\).
  Complex-bilinear expansion against its conjugate produces the four real
  sectional terms and two mixed terms.  The first Bianchi identity combines
  the mixed terms into \(2R(y,x,v,u)\), giving the displayed formula.
\end{proof}

\begin{remark}[special cases of the complex sectional curvature]
  \label{rem:pet-ch9-complex-sectional-curvature-special-cases}
  \lean{PetersenLib.complexSectionalCurvature_specialCases}
  \source{petersen:page-0371}
  \uses{lem:pet-ch9-complex-sectional-curvature-real-formula}
  Three specializations of the preceding formula are used:
  \begin{enumerate}
    \item $y=v=0$ gives the ordinary sectional quantity $R(x,u,u,x)$;
    \item orthonormal $x,y,u,v$ gives the \emph{isotropic curvature};
    \item $u=v$ gives $2R(x,u,u,x)+2R(y,u,u,y)$, called the
      \emph{second Ricci curvature} when $x,y,u$ are orthonormal.
  \end{enumerate}
\end{remark}

\begin{proposition}[normal tensors and complex sectional curvature]
  \label{prop:pet-ch9-normal-tensor-nonneg-complex-sectional}
  \lean{PetersenLib.normalTensor_nonneg_of_complexSectionalCurvature_nonneg}
  \source{petersen:page-0371}
  \uses{prop:pet-ch9-02tensor-hat-formula, def:pet-ch9-complex-sectional-curvature}
  \dcref{ch9:9.4.12}
  If the endomorphism $H$ associated with a $(0,2)$-tensor $h$ is normal, then
  nonnegative complex sectional curvature implies
  $g(\mathfrak R(\hat h),\hat h)\ge0$.
\end{proposition}
\begin{proof}
  \uses{prop:pet-ch9-normal-tensor-nonneg-complex-sectional}
  \sketch
  Choose a complex orthonormal eigenbasis \(H(e_i)=\lambda_ie_i\).  By
  \cref{prop:pet-ch9-02tensor-hat-formula},
  \[
    g(\mathfrak R(\hat h),\hat h)
      =\sum_{i,j}|\lambda_i-\bar\lambda_j|^2
        g(\mathfrak R(e_i\wedge e_j),\bar e_i\wedge\bar e_j).
  \]
  Every summand is nonnegative under the stated hypothesis.
\end{proof}

\begin{proposition}[general tensors and complex sectional curvature]
  \label{prop:pet-ch9-general-tensor-nonneg-complex-sectional}
  \lean{PetersenLib.tensor_nonneg_of_complexSectionalCurvature_nonneg}
  \source{petersen:page-0372}
  \uses{prop:pet-ch9-normal-tensor-nonneg-complex-sectional}
  \dcref{ch9:9.4.13}
  At a point with nonnegative complex sectional curvature,
  \[
    g(\mathfrak R(\hat h),\hat h)\ge0
  \]
  for every $(0,2)$-tensor $h$.
\end{proposition}
\begin{proof}
  \uses{prop:pet-ch9-general-tensor-nonneg-complex-sectional}
  \sketch
  Decompose \(h=h_s+h_a\) into symmetric and alternating parts.  Both
  associated operators are normal.  Symmetry of \(\hat h_s(e_i,e_j)\) and
  antisymmetry of \(\hat h_a(e_i,e_j)\) make the mixed contraction vanish, so
  \[
    g(\mathfrak R(\hat h),\hat h)
      =g(\mathfrak R(\hat h_s),\hat h_s)
       +g(\mathfrak R(\hat h_a),\hat h_a).
  \]
  Each term is nonnegative by
  \cref{prop:pet-ch9-normal-tensor-nonneg-complex-sectional}.
\end{proof}

\section{Exercises}

\begin{remark}[splitting of the universal cover]
  \label{rem:pet-ch9-exercise-9-6-1}
  \lean{PetersenLib.exercise9_6_1}
  \source{petersen:page-0373}
  \uses{cor:pet-ch9-betti1-bound-flat-torus, thm:pet-ch9-bochner-vanishing-1forms}
  \dcref{ch9:9.6.1}
  Let compact $(M^n,g)$ have $b_1(M)=k$ and $\operatorname{Ric}\ge0$.  Prove
  the product decomposition
  \[
    (\widetilde M,g)=(N,h)\times(\mathbb R^k,g_{\mathbb R^k}).
  \]
  Also construct a case with $b_1<n$ but
  $(\widetilde M,g)=(\mathbb R^n,g_{\mathbb R^n})$.
\end{remark}

\begin{remark}[frame identity for a symmetric derivative]
  \label{rem:pet-ch9-exercise-9-6-2}
  \lean{PetersenLib.exercise9_6_2}
  \source{petersen:page-0373}
  \uses{prop:pet-ch9-vector-field-1form-duality}
  \dcref{ch9:9.6.2}
  For an arbitrary orthonormal frame $E_i$, derive
  \[
    \sum_i g(\nabla_{E_i}X,\nabla_XE_i)=0
  \]
  directly from symmetry of $v\mapsto\nabla_vX$; no parallelism of $E_i$
  along $X$ may be assumed.
\end{remark}

\begin{remark}[orthogonality of images and kernels]
  \label{rem:pet-ch9-exercise-9-6-3}
  \lean{PetersenLib.exercise9_6_3}
  \source{petersen:page-0373}
  \uses{def:pet-ch9-codifferential}
  \dcref{ch9:9.6.3}
  In $\Omega^k(M)$ for an oriented Riemannian manifold, establish
  \[
    (\operatorname{im}d^{k-1})^\perp=\ker\delta^{k-1},
    \qquad
    \operatorname{im}d^{k-1}\subset(\ker\delta^{k-1})^\perp.
  \]
\end{remark}

\begin{remark}[elliptic regularity for the connection Laplacian]
  \label{rem:pet-ch9-exercise-9-6-4}
  \lean{PetersenLib.exercise9_6_4}
  \source{petersen:page-0373}
  \uses{def:pet-ch9-lichnerowicz-laplacian}
  \dcref{ch9:9.6.4}
  For compact $(M^n,g)$ and a tensor bundle $E$, let $W^{k,2}(E)$ be the
  completion of $\Gamma(E)$ in the norm
  $\sum_{i=0}^k\|\nabla^iT\|_2^2$.  Thus $W^{1,2}$ has a weak derivative with
  $(\nabla T,\nabla S)=(T,\nabla^*\nabla S)$ for smooth $S$.  Smooth sections
  are dense; Sobolev estimates identify $\bigcap_kW^{k,2}(E)$ with
  $\bigcap_kW^{k,p}(E)$ for every $p$, and the section 7.1.5 regularity result
  makes $W^{1,p}$ sections H\"older for $p>n$.  Consequently
  $\Gamma(E)=\bigcap_kW^{k,2}(E)$.
  \begin{enumerate}
    \item[(1)] Establish, for each smooth $T$, the commutator formula
      $\nabla^*\nabla(\nabla^kT)-\nabla^k(\nabla^*\nabla T)
      =\sum_{i=0}^kC_i^k(\nabla^{k-1}R\otimes\nabla^iT)$ with suitable
      contractions $C_i^k$.
    \item[(2)] Suppose $T\in W^{1,2}(E)$ and $T'\in W^{k,2}(E)$ satisfy
      $(T',S)=(\nabla T,\nabla S)$ for every smooth $S$, i.e.
      $\nabla^*\nabla T=T'$ weakly.  Prove $T\in W^{k+1,2}(E)$ by defining
      weak derivatives inductively through
      $(\nabla^{l+1}T,\nabla^{l+1}S)=(\nabla^lT',\nabla^lS)
      +\sum_{i=0}^l(C_i^l(\nabla^{l-i}R\otimes\nabla^iT),\nabla^lS)$.
    \item[(3)] Deduce smoothness of $T$ and the strong equation
      $\nabla^*\nabla T=T'$ whenever $T'$ is smooth.
  \end{enumerate}
\end{remark}

\begin{remark}[spectral theorem for a Lichnerowicz Laplacian]
  \label{rem:pet-ch9-exercise-9-6-5}
  \lean{PetersenLib.exercise9_6_5}
  \source{petersen:page-0373,petersen:page-0374}
  \uses{def:pet-ch9-lichnerowicz-laplacian, rem:pet-ch9-exercise-9-6-4, cor:pet-ch9-ric-nonneg-curvature-operator}
  \dcref{ch9:9.6.5}
  Let compact $(M^n,g)$ satisfy $\operatorname{diam}M\le D$ and
  $\operatorname{Ric}\ge k$, and let the rank-$m$ tensor bundle $E$ carry
  $\Delta_L$.  Prove the spectral theorem and the orthogonal splitting
  $\Gamma(E)=\ker\Delta_L\oplus\operatorname{im}\Delta_L$ by establishing:
  \begin{enumerate}
    \item[(1)] On $W^{1,2}(E)$, prove that
      $(\Delta_LT,S)=(\nabla T,\nabla S)+c(\operatorname{Ric}(T),S)$ is
      well-defined and symmetric.
    \item[(2)] For $k\le0$, prove
      $\inf_{T\in W^{1,2},\|T\|_2=1}(\Delta_LT,T)>cCk$, with $C$ from
      Corollary 9.3.4.
    \item[(3)] Show that every unit sequence with bounded
      $(\Delta_LT_i,T_i)$ has an $L^2$-convergent subsequence, weakly convergent
      in $W^{1,2}(E)$ by theorem 7.1.18.
    \item[(4)] For a closed $\Delta_L$-invariant
      $V\subset W^{1,2}(E)$, show the Rayleigh infimum
      $\lambda=\inf_{T\in V,\|T\|_2=1}(\Delta_LT,T)$ is attained, and use
      parts (1)--(3) of Exercise 9.6.4 to obtain
      $T\in\Gamma(E)$ with $\Delta_LT=\lambda T$.
    \item[(5)] If finite-dimensional $V\subset\Gamma(E)$ is spanned by
      eigentensors $\Delta_LT_i=\lambda_iT_i$, prove
      $\dim V\le mC(n,\max_i\lambda_i,c,D^2k)$.
    \item[(6)] Prove finite dimensionality of every eigenspace and discreteness
      of the spectrum $\lambda_1<\lambda_2<\cdots\to\infty$.
    \item[(7)] Prove orthogonality of distinct eigenspaces and density of their
      direct sum in $\Gamma(E)$.
    \item[(8)] Conclude $\Gamma(E)=\ker\Delta_L\oplus\operatorname{im}\Delta_L$.
  \end{enumerate}
\end{remark}

\begin{remark}[Euclidean rigidity outside a compact set]
  \label{rem:pet-ch9-exercise-9-6-6}
  \lean{PetersenLib.exercise9_6_6}
  \source{petersen:page-0374}
  \uses{thm:pet-ch9-bochner-vanishing-1forms}
  \dcref{ch9:9.6.6}
  Suppose $M^n-K\cong\mathbb R^n-C$ isometrically for compact
  $K\subset M$ and $C\subset\mathbb R^n$.  Prove that
  $\operatorname{Ric}_g\ge0$ forces $(M,g)$ to be Euclidean.  One may either
  embed a neighborhood of $K$ into a metric on $T^n$ that is flat elsewhere,
  or extend parallel one-forms from $\mathbb R^n-C$ harmonically and apply
  Bochner vanishing.
\end{remark}

\begin{remark}[harmonic one-forms on Einstein manifolds]
  \label{rem:pet-ch9-exercise-9-6-7}
  \lean{PetersenLib.exercise9_6_7}
  \source{petersen:page-0374}
  \uses{def:pet-ch9-lichnerowicz-laplacian, thm:pet-ch9-weitzenbock-hodge-laplacian}
  \dcref{ch9:9.6.7}
  Prove that on an Einstein manifold every harmonic one-form is an eigenform
  of $\nabla^*\nabla$.
\end{remark}

\begin{remark}[mixed Bochner formulas]
  \label{rem:pet-ch9-exercise-9-6-8}
  \lean{PetersenLib.exercise9_6_8}
  \source{petersen:page-0374}
  \uses{prop:pet-ch9-bochner-formula-1forms}
  \dcref{ch9:9.6.8}
  If both $\nabla X$ and $\nabla Y$ are symmetric, derive the polarized
  Bochner identities for
  $\operatorname{Hess}(g(X,Y)/2)$ and $\Delta(g(X,Y)/2)$.
\end{remark}

\begin{remark}[tensor inner-product identities]
  \label{rem:pet-ch9-exercise-9-6-9}
  \lean{PetersenLib.exercise9_6_9}
  \source{petersen:page-0374,petersen:page-0375}
  \uses{prop:pet-ch9-vanishing-connection-laplacian}
  \dcref{ch9:9.6.9}
  For same-type tensors $s_1,s_2$, polarize the identity
  $\Delta\tfrac12|s|^2=|\nabla s|^2-g(\nabla^*\nabla s,s)$ to prove
  \[
    \Delta g(s_1,s_2)=2g(\nabla s_1,\nabla s_2)+g(\nabla^*\nabla s_1,s_2)
    +g(s_1,\nabla^*\nabla s_2).
  \]
  Apply it to forms to obtain Bochner identities for their inner products.
  Next set $\omega(v)=g(\nabla_vs_1,s_2)$, which represents half the
  differential of $g(s_1,s_2)$, and verify
  \[
    -\delta\omega=g(\nabla s_1,\nabla s_2)-g(\nabla^*\nabla s_1,s_2),
  \]
  \[
    d\omega(X,Y)=g(R(X,Y)s_1,s_2)-g(\nabla_Xs_1,\nabla_Ys_2)+g(\nabla_Ys_1,\nabla_Xs_2).
  \]
\end{remark}

\begin{remark}[skew endomorphisms acting on volume and curvature]
  \label{rem:pet-ch9-exercise-9-6-10}
  \lean{PetersenLib.exercise9_6_10}
  \source{petersen:page-0375}
  \uses{conv:pet-ch9-so-metric-convention}
  \dcref{ch9:9.6.10}
  On an $n$-manifold, verify the following actions of an endomorphism $L$:
  \begin{enumerate}
    \item[(1)] If $L$ is skew-symmetric, then $L\omega=0$ for every $n$-form.
    \item[(2)] In dimension two, every skew-symmetric $L$ satisfies $LR=0$.
    \item[(3)] Without a skew-symmetry assumption,
      $L\,\mathrm{vol}=\operatorname{tr}(L)\,\mathrm{vol}$.
  \end{enumerate}
\end{remark}

\begin{remark}[Simons' estimates for Codazzi tensors]
  \label{rem:pet-ch9-exercise-9-6-11}
  \lean{PetersenLib.exercise9_6_11}
  \source{petersen:page-0375}
  \uses{def:pet-ch9-codazzi-harmonic-tensor, thm:pet-ch9-symmetric-tensor-lichnerowicz-identity, prop:pet-ch9-vanishing-connection-laplacian}
  \dcref{ch9:9.6.11}
  Let $h$ be a constant-trace symmetric Codazzi tensor on compact $(M,g)$.
  \begin{enumerate}
    \item[(1)] Deduce $\nabla h=0$ from $\sec\ge0$.
    \item[(2)] Under $\sec>0$, deduce $h=cg$ for a constant $c$.
    \item[(3)] Assuming $\operatorname{tr}h=0$ and the Gauss equation
      $R(X,Y,Z,W)=h(X,W)h(Y,Z)-h(X,Z)h(Y,W)$, prove
      $\Delta\tfrac12|h|^2\ge|\nabla h|^2-|h|^4$.
  \end{enumerate}
\end{remark}

\begin{remark}[constant-mean-curvature hypersphere rigidity]
  \label{rem:pet-ch9-exercise-9-6-12}
  \lean{PetersenLib.exercise9_6_12}
  \source{petersen:page-0375,petersen:page-0376}
  \uses{def:pet-ch9-codazzi-exterior-derivative, rem:pet-ch9-exercise-9-6-11}
  \dcref{ch9:9.6.12}
  For an isometric immersion $(M^n,g)\hookrightarrow\mathbb R^{n+1}$:
  \begin{enumerate}
    \item[(1)] prove that the second fundamental form $\mathrm{II}$ is
      Codazzi;
    \item[(2)] prove Liebmann's theorem: constant mean curvature and
      $\mathrm{II}\ge0$ imply that $(M,g)$ is a constant-curvature sphere.
      The nonnegativity assumption cannot simply be omitted, since Wente's
      immersed tori have constant mean curvature.
  \end{enumerate}
\end{remark}

\begin{remark}[Berger's parallel-Ricci criterion]
  \label{rem:pet-ch9-exercise-9-6-13}
  \lean{PetersenLib.exercise9_6_13}
  \source{petersen:page-0376}
  \uses{cor:pet-ch9-ricci-harmonic-iff-curvature-harmonic, thm:pet-ch9-symmetric-tensor-lichnerowicz-identity}
  \dcref{ch9:9.6.13}
  Prove that $\nabla\operatorname{Ric}=0$ on every compact manifold satisfying
  $\nabla^*R=0$ and $\sec\ge0$.
\end{remark}

\begin{remark}[isometries and harmonic forms]
  \label{rem:pet-ch9-exercise-9-6-14}
  \lean{PetersenLib.exercise9_6_14}
  \source{petersen:page-0376}
  \uses{thm:pet-ch9-hodge-theorem, def:pet-ch9-hodge-cohomology}
  \dcref{ch9:9.6.14}
  Fix a Killing field $K$ and a harmonic form $\omega$ on closed $(M,g)$.
  \begin{enumerate}
    \item[(1)] Prove $\mathcal L_K\omega=0$ by first showing that this Lie
      derivative is harmonic.
    \item[(2)] Prove $i_K\omega$ is closed, and show that it need not be
      harmonic.
    \item[(3)] Deduce invariance of every harmonic form under
      $\operatorname{Iso}_0(M)$.
    \item[(4)] Exhibit failure of invariance under the full group
      $\operatorname{Iso}(M)$.
  \end{enumerate}
\end{remark}

\begin{remark}[nonexact powers of the K\"ahler form]
  \label{rem:pet-ch9-exercise-9-6-15}
  \lean{PetersenLib.exercise9_6_15}
  \source{petersen:page-0376}
  \uses{rem:pet-ch9-poincare-duality}
  \dcref{ch9:9.6.15}
  Let $\omega$ be the parallel nondegenerate K\"ahler form of a closed
  K\"ahler manifold.  Prove that each wedge power
  $\omega^k=\omega\wedge\cdots\wedge\omega$ is closed and nonexact.  Use that
  $\omega^{\dim M/2}$ is proportional to the volume form, and conclude that
  every even homology group is nonzero.
\end{remark}

\begin{remark}[vector-bundle-valued differential forms]
  \label{rem:pet-ch9-exercise-9-6-16}
  \lean{PetersenLib.exercise9_6_16}
  \source{petersen:page-0376}
  \uses{def:pet-ch9-codazzi-exterior-derivative, def:pet-ch9-weitzenbock-curvature-operator}
  \dcref{ch9:9.6.16}
  For a tensor bundle $E\to M$, set $\Omega^0(M,E)=\Gamma(E)$ and let
  $\Omega^p(M,E)$ be the alternating $E$-valued $p$-forms.  Establish:
  \begin{enumerate}
    \item[(1)] the natural left and right wedge actions of $\Omega^*(M)$ on
      $\Omega^*(M,E)$;
    \item[(2)] the wedge product
      $\Omega^p(M,\operatorname{Hom}(E,E))\times\Omega^q(M,E)\to\Omega^{p+q}(M,E)$.
    \item[(3)] a connection-dependent differential
      $d^\nabla:\Omega^p(M,E)\to\Omega^{p+1}(M,E)$ satisfying the Leibniz
      rule and $d^\nabla s=\nabla s$ for $s\in\Gamma(E)$;
    \item[(4)] with $R(X,Y)s$ viewed as an
      $\operatorname{Hom}(E,E)$-valued two-form, the identities
      $(d^\nabla)^2s=R\wedge s$ for $s\in\Omega^p(M,E)$ and
      $d^\nabla R=0$ (the second Bianchi identity).
  \end{enumerate}
\end{remark}

\begin{remark}[Codazzi tensors with bundle values]
  \label{rem:pet-ch9-exercise-9-6-17}
  \lean{PetersenLib.exercise9_6_17}
  \source{petersen:page-0377}
  \uses{rem:pet-ch9-exercise-9-6-16, def:pet-ch9-codazzi-harmonic-tensor}
  \dcref{ch9:9.6.17}
  Specialize Exercise 9.6.16 to $E=TM$, so that
  $\Omega^1(M,TM)=\operatorname{Hom}(TM,TM)$.
  \begin{enumerate}
    \item[(1)] prove $d^\nabla s=0$ exactly for Codazzi tensors;
    \item[(2)] develop, for $d^\nabla$-closed $s\in\Omega^p(M,E)$, a
      Bochner identity and a parallelism theorem under divergence-free and
      curvature-nonnegative hypotheses;
    \item[(3)] show $\nabla X$ is Codazzi iff
      $R(\cdot,\cdot)X=0$, give a nonparallel example, and determine whether
      an analogous formula applies to the exact tensor $\nabla X=d^\nabla X$
      without closedness.
  \end{enumerate}
\end{remark}

\begin{remark}[Thomas' Gauss-to-Codazzi implication]
  \label{rem:pet-ch9-exercise-9-6-18}
  \lean{PetersenLib.exercise9_6_18}
  \source{petersen:page-0377}
  \uses{rem:pet-ch9-exercise-9-6-17, def:pet-ch9-codazzi-harmonic-tensor}
  \dcref{ch9:9.6.18}
  In dimension $n\ge3$, prove that the Gauss relation
  $\mathfrak R=S\wedge S$ forces $d^\nabla S=0$ when $\det S\ne0$.  Use the
  second Bianchi identity and analyze
  \[
    T\longmapsto T\wedge S:
    \operatorname{Hom}(\Lambda^2V,V)\to\operatorname{Hom}(\Lambda^3V,\Lambda^2V).
  \]
  This map is injective only for $\operatorname{rank}S\ge4$.
\end{remark}
